\subsection{From certified final-feature input errors to covariance-score signs}
\label{sec:p4-final-feature-chain}
The covariance-coordinate theorem previously assumed
an externally supplied matrix-level error radius.
We now derive that radius \emph{deterministically}
from a more elementary, equally conditional assumption:
the \emph{final stored feature matrix} $\widehat X\in
\mathbb Q^{n\times p}$ is related to the hypothetical
unrounded feature matrix $X$ by entrywise
$|X_{ij}-\widehat X_{ij}|\leq\eta$.
Here $\eta$ is an \emph{externally justified bound
on the final features}, \emph{not} a property we
infer from an EDF's raw 16-bit conversion.
No non-linear EEG preprocessing, tangent transform
or numerical matrix decomposition is covered by
this assumption without an additional verified
forward-error analysis.

For two classes, let $H$ be the symmetric
block-diagonal within-class centring projection,
$R=HX$, $\widehat R=H\widehat X$ and
$\rho I$ a fixed identical positive ridge.
Then the pooled within-class covariances are
\[
S=\frac{R^\top R}{n-2}+\rho I,\qquad
\widehat S=\frac{\widehat R^\top\widehat R}{n-2}+\rho I.
\]
Writing $E=X-\widehat X$ gives
$\|HE\|_2\leq\|E\|_F\leq\eta\sqrt{np}$.
By expanding the Gram product,
\[
\|S-\widehat S\|_2\leq
\frac{2\|\widehat R\|_F\eta\sqrt{np}
       +np\eta^2}{n-2}
\equiv\epsilon_S.
\]
For the difference $d$ of two empirical class
means and any positive-minus-negative example
difference $v$, the same entrywise assumption
implies
\[
\|d-\widehat d\|_2\leq 2\eta\sqrt{p}
\equiv\epsilon_d,\qquad
\|v-\widehat v\|_2\leq 2\eta\sqrt{p}
\equiv\epsilon_v.
\]
Replacing all square roots with \emph{upward exact-rational}
approximations gives a machine-checkable rational
outer bound, which is submitted to the
covariance-coordinate resolvent criterion in
Section~\ref{sec:p4-cov-resolvent}.
If it returns a strict sign, that sign holds for
\emph{every} exact feature matrix in the specified
entrywise input box and every $\lambda\in[0,1]$.
No independence assumptions on feature errors,
no fixed eigenvectors and no sampled noise law
are required for this conditional conclusion.

The included synthetic exact-rational test uses
64 seeded two-class, 24-record final-feature matrices
with a common positive ridge. Of these deliberately
well-separated fixtures, 48 have proved positive
score signs and 16 proved negative signs over the
entire continuous $\lambda$ interval and every
admissible entrywise feature perturbation; no
candidate was inconclusive in this \emph{easy}
synthetic subset. Independently reconstructed
admissible feature perturbations produced 576 exact
rational score checks without a contradiction.
A separate, deliberately adversarial \emph{exact
counterexample} perturbs one positive and one
negative record, each by $\eta=0.1$: the nominal
score is positive and an admissible score is
negative across tested shrinkage values; the
algorithm correctly returns \texttt{UNKNOWN}.
This demonstrates that the apparatus can
refuse certification when input uncertainty
permits a genuine decision reversal.
The program recomputes all class means,
within-class residuals, covariance and pair
differences rather than merely injecting
a matrix-level error. The successful code,
exact samples and fail-closed counterexample
are archived in
\url{https://github.com/lindicaphxag-tech/kaggle/actions/runs/38389194383}.
These are conditional operator regressions,
not real EEG classification, a patient-level
risk model, or proof that any unverified
preprocessing/tangent pipeline satisfies
the assumed final-feature $\eta$ bound.
